\section{From tagged permutations to amortized simulation}
\label{sec:descent}

Fix a finite field $F$ with largest spectral flat dimension $k$. A
\emph{tick} is a trivalent bit network and its \emph{pins} are its open
data bits. For a tick of size at most $V$, we can reduce the size per
copy to $V/(\log_2 V)^A$ for any fixed $A$, provided the tick has
sufficiently few pins. Each copy retains its own pins. The approximation
must work for every joint choice of their values.

\begin{theorem}[Amortized simulation]
\label{thm:collapse}
For every sufficiently large fixed integer $A$ there is a constant
integer $A'>A$ with the following property. Let $V=2^w$. For all
sufficiently large integers $w$, let a tick have size at most $V$ and
at most $V/w^{A'}$ pins. For some positive integer $q\le V^{o(1)}$,
there is a trivalent network of size at most $qV/w^A$ approximating
the product of $q$ copies of the tick, with gap at least $V^{1-o(1)}$
pointwise in their separate pins. Both asymptotic envelopes are
uniform in the tick.
\end{theorem}

Two costs determine whether the saving can be iterated.
Proposition~\ref{prop:oracle-descent} reduces the size per tick by a
power of the address length, using exact tagged-permutation gates.
Theorem~\ref{thm:cheaptag} replaces those gates at a smaller overhead,
but each replacement exposes new catalyst bits. At the next descent
step every exposed bit counts toward the compiler's boundary cost. We
control that cost by placing copies on a grid and passing each
catalyst wire successively through the cells of a coordinate line.
Only the two ends of the line remain exposed. The price estimate below
quantifies how much of the original gap this operation preserves.

\subsection{Pricing error on a grid}

Consider the grid $\{3,6,\ldots,3g\}^k$. For direction $j$, place
two \emph{foci} at $\pm Re_j$, with $R\ge8g$. There are $E$ external
pairs of bits, each assigned a direction and having one end at each
focus of that direction. In a fixed wiring, choose a grid line in the
assigned direction for each pair. Its \emph{rail} runs from the
negative focus to the first cell on the line, visits its cells in order
and ends at the positive focus.

At each cell the exact wiring uses one factor $f\otimes I_X$. Attach
each of its catalyst pairs to a distinct rail passage in the prescribed
direction. The minus and plus catalyst ports form the incoming and
outgoing ends of that passage. All other passages and the data ports
may participate in an arbitrary local contraction. The rails supply all
intercell connections, and their endpoints at the foci are the only
outputs. The approximate wiring replaces the factor at every cell by
$X$ with gap at least $D>0$. Its data are exposed except for pointwise
pins, whose values must be fixed. We allow both constructions to be
summed over the same collection of at most $2^\ell$ choices of wiring,
local factors and settings.

\begin{lemma}[Lens estimate]\label{lem:lens}
The difference of the two sums has price at most
\[
  2RE-D+O_k(Eg^2/R)+2kR(\ell+g^k).
\]
More precisely, putting $L=g^k$ and taking $c$ to be the binary diagonal
price on the set of all $2k$ foci, it has price at most
\begin{equation}\label{eq:lens-binomial}
  2RE+C_k Eg^2/R+c\ell+
       c\log_2\big((1+2^{-D/c})^L-1\big)
\end{equation}
for a constant $C_k$ depending only on $k$.
\end{lemma}

\begin{proof}
Expand the difference for one wiring by multilinearity. Each term has
a nonempty set of cells at which it uses $X-f\otimes I_X$. At cell $x$,
place the incoming and outgoing ends of a direction-$j$ passage at
$x-e_j$ and $x+e_j$, and the data operations at $x$. Allow two units of
price per passage. This pays for the through identity on a catalyst pair
or, on another passage, for moving both ends to $x$. A cell error saves $D$ from the total allowance at that
cell by \eqref{eq:catalyst-gap} in the translated frame. An interior link
along a rail has length $3-2=1$. Each of its two end links has a longitudinal
component $a\ge R/2$ and perpendicular displacement $u$ with
$\|u\|^2\le 9(k-1)g^2$. Its excess length is
\[
   \sqrt{a^2+\|u\|^2}-a
      =\frac{\|u\|^2}{\sqrt{a^2+\|u\|^2}+a}
      \le \frac{\|u\|^2}{2a}.
\]
The longitudinal parts telescope, including the passage allowances, to $2R$
per rail (Figure~\ref{fig:spectral-lens}). A summand with $i$ erroneous
cells therefore has price at most $2RE+C_kEg^2/R-iD$.

There are at most $2^\ell\binom{L}{i}$ terms of this kind.
Equation~\eqref{eq:price-ordinary-sum} gives \eqref{eq:lens-binomial} by
summing these spectral values. Simply bounding the number of terms by
$2^{\ell+L}$ and using $i\ge1$, with $c\le2kR$, gives the first assertion.
The same output frame is used for every wiring, even when each selector sends
the external pairs to different grid lines.
\end{proof}

\begin{figure}[tbp]
\centering
\begin{tikzpicture}[scale=.9,line cap=round,line join=round,font=\small]
  \definecolor{lensfirst}{RGB}{27,91,130}
  \definecolor{lenssecond}{RGB}{165,72,32}
  \draw[gray!38,dashed] (-3.55,0)--(3.55,0);
  \draw[gray!38,dashed] (0,-3.55)--(0,3.55);
  \foreach \v in {.44,.88,1.32}{
     \draw[gray!40] (.44,\v)--(1.32,\v);
     \draw[gray!40] (\v,.44)--(\v,1.32);
  }
  \draw[lensfirst!55,densely dotted] (1.32,.88)--(1.32,0)
       node[midway,right,font=\scriptsize,text=black] {$\|u\|$};
  \node[below,font=\scriptsize] at (2.43,0) {$a$};
  \draw[lensfirst,thick] (-3.55,0)--(.44,.88)--(1.32,.88)--(3.55,0);
  \draw[lenssecond,thick,dashed] (0,-3.55)--(.88,.44)--(.88,1.32)--(0,3.55);
  \foreach \x in {.44,.88,1.32}{
    \foreach \y in {.44,.88,1.32}{
      \fill[black] (\x,\y) circle (1.45pt);
    }
  }
  \draw[black,fill=white] (.88,.88) circle (2.5pt);
  \fill[lensfirst] (-3.55,0) circle (2pt) node[left] {$-Re_1$};
  \fill[lensfirst] (3.55,0) circle (2pt) node[right] {$Re_1$};
  \fill[lenssecond] (0,-3.55) circle (2pt) node[right] {$-Re_2$};
  \fill[lenssecond] (0,3.55) circle (2pt) node[right] {$Re_2$};
  \begin{scope}[shift={(7.65,.3)}]
    \draw[black!45,rounded corners] (-2.78,-2.12) rectangle (2.78,2.02);
    \draw[lensfirst,thick] (-1.18,0)--(1.18,0);
    \draw[lenssecond,thick,dashed] (0,-1.18)--(0,1.18);
    \fill[lensfirst] (-1.18,0) circle (2pt) node[left] {$x-e_1$};
    \fill[lensfirst] (1.18,0) circle (2pt) node[right] {$x+e_1$};
    \fill[lenssecond] (0,1.18) circle (2pt) node[above] {$x+e_2$};
    \fill[lenssecond] (0,-1.18) circle (2pt) node[below] {$x-e_2$};
    \draw[fill=white] (0,0) circle (2.5pt) node[above right] {$x$};
    \node at (0,-2.56) {$H(X-f\otimes I_X)\le 2(b_1+b_2)-D$};
  \end{scope}
\end{tikzpicture}
\caption{Two rails on a planar grid. At the circled cell each rail uses
the ordered pair of sites shown at right. Its projection from focus to focus
has length $2R$. Only the sloping end links add excess length. The marked
link has length $\sqrt{a^2+\|u\|^2}$ and contributes at most
$\|u\|^2/(2a)$ beyond its projection.}
\label{fig:spectral-lens}
\end{figure}

Splicing a line of $g$ copies reduces its number of external catalyst
pairs by a factor $g$. To transfer this construction back to unit frames,
choose the foci far enough away that the total excess rail length is
small compared with the gap.

\begin{lemma}[Grid reuse]\label{lem:reuse}
Suppose a scalar network at fixed pins has gap at least $D=V^{1-o(1)}$ and
uses $B\le V^{1+o(1)}$ catalyst pairs. Put $g=V^{o(1)}$, a positive
integer. Take $g^k$ copies with separate pins and splice consecutive
catalyst ports on each line in their common direction. The resulting
network approximates the product pointwise with gap $V^{1-o(1)}$ and has
$g^{k-1}B$ external catalyst pairs.
\end{lemma}
\begin{proof}
In each direction there are $g^{k-1}$ grid lines. After fixing the pins
the target is the product of the scalar values times the identity on the
$E=g^{k-1}B$ remaining pairs. Take $D\le V$ by decreasing the claimed
gap if necessary. With $\ell=0$, choose $R=V^{o(1)}$, $R\ge8g$, so that
$Eg^2/R=o(D)$. This is possible since $Eg^2/D\le V^{o(1)}$, and such
an $R$ still gives $Rg^k=o(D)$. For any prescribed unit output frame with
placement $i:\mathbb R^k\to\mathcal X$, use the spectral space with prices
and distances multiplied by $R$ and the isometry $u\mapsto i(u/R)$ into
that rescaled space. The lens foci are now exactly $i(\pm e_j)$.
Lemma~\ref{lem:lens} gives an error price at most $2RE-D/2$ there.
Dividing by $R$ returns to the prescribed frame with gap at least
$D/(2R)=V^{1-o(1)}$. Bounds uniform in the original networks can use a
common choice of $R$.
\end{proof}

\subsection{Iteration of the descent step}

\begin{proof}[Proof of Theorem~\ref{thm:collapse}]
Take $a\in(0,1)$ from Proposition~\ref{prop:oracle-descent}. Apply
Theorem~\ref{thm:cheaptag} with $r=a/4$, so its network for
$\operatorname{Tag}(h,h)$ has size $O(2^h h^{1+r})$ and exposed gap
at least $2^{h-C\sqrt h}$ for a fixed $C$. Choose
\[
 b=a/2,\qquad L=\lceil A/b\rceil,\qquad
 A'=\lceil bL+4\rceil,\qquad
 h=w-\lceil A'\log_2 w\rceil,\qquad N=2^h .
\]
Thus $h\sim w$ and $V/(2w^{A'})\le N\le V/w^{A'}$.

After $i$ steps we shall have a network $G_i$ representing $q_i$ ticks
pointwise with gap $V^{1-o(1)}$ or better. Its size per tick is at most
$V/w^{bi}$, and its total number of open bits per tick, including the
two ends of every catalyst pair, is at most $2V/w^{A'}$.
We also require $q_i\le V^{o(1)}$. The tick itself satisfies these
conditions at $i=0$ with $q_0=1$ and zero error.

Suppose $i<L$. First compile the product using exact tagged gates.
Normalize $G_i$, write $e$ for its open bit count, and
let $m$ be the larger of its normalized size and $N$. Then
$m\le N^2$ for large $w$, since $q_i V=V^{1+o(1)}$ and
$N=V^{1-o(1)}$. Proposition~\ref{prop:oracle-descent} gives an exact
network for $G_i^{\otimes d}$ with
$d=2^{O((\log h)^2)}$. Its size is $V^{1+o(1)}$. The number of tag
calls is $V^{o(1)}$, since each has $3Nh$ bit incidences and each
wire contributes at most two incidences to this count. Replace every
call by the approximation above, retaining all of its catalyst ports.

First apply Lemma~\ref{lem:substitution} to $d$ copies of $G_i$ in
the product of $dq_i$ ticks, with the original pins fixed.
The resulting error bound depends only on the tensor, hence applies
to the exact oracle realization. Apply the lemma a second time to its
tag calls, with their data exposed so that the internal indices may be
contracted. The two applications each lose at most $V^{o(1)}$ in
absolute gap to summation over error terms. The tag gaps
$N2^{-C\sqrt h}$ and the inherited gap are at least $V^{1-o(1)}$.
Thus the result still has gap at least $V^{1-o(1)}$ pointwise in all
$dq_i$ copies' pins.

The compiler charges every bit incident to a tag call. Replacing the
calls therefore increases size by a factor of at most $O(h^r)$.
The compiler's interior cost and boundary cost give, respectively, the
two terms in the following bound on size per tick.
\[
 O\left(\frac{V}{w^{bi+a-r}}+
         \frac{V\operatorname{polylog} w}{w^{A'-r}}\right)
              \le \frac{V}{w^{b(i+1)}}
\]
for sufficiently large $w$. Indeed normalization costs a constant
factor, the floor $m\ge N$ contributes at most the second term, and
$e/q_i\le 2V/w^{A'}$. The two exponents have slack because
$a-r>b$ and $A'-r>bL$.

The new catalyst ports can outnumber the original pins. Apply
Lemma~\ref{lem:reuse} to a grid of side $g=\lceil w^{A'}\rceil$.
Replication by $g^k$ and splicing along the lines in each catalyst
direction preserves the pointwise gap $V^{1-o(1)}$. Indeed the number
of pairs before reuse is at most the total network size
$V^{1+o(1)}$, and $g=V^{o(1)}$. Each line has $g$ copies and keeps
one pair of ends in place of $g$ such pairs. Hence the number of open
catalyst bits per tick falls by a factor $g$, and is now at most
$V/g$. The number of pins per tick is still at most $V/w^{A'}$.
Splicing does not increase the size per tick. This proves the inductive
bounds with $q_{i+1}=g^k d q_i\le V^{o(1)}$.

After the fixed number $L$ of steps, $bL\ge A$ gives the asserted size.
Normalization, the compiler, the tag estimates, and grid reuse allow
uniform bounds at each step. Taking their composition a fixed number
of times gives the stated envelopes.
\end{proof}
